\subsection{Needspace -- Page Break Control}

The \texttt{needspace} package starts a new page when there is not
enough space for the material that follows. Load it in the preamble:

\begin{lstlisting}[style=latex]
\usepackage{needspace}
\end{lstlisting}

\subsubsection{Basic Usage}

Place \verb|\Needspace| between paragraphs, immediately before
the content that requires a minimum amount of space:

\begin{lstlisting}[style=latex]
Previous paragraph.

\Needspace{5\baselineskip}
This paragraph starts on the next page if less than
five baseline spacings remain on the current page.
\end{lstlisting}

If enough space remains, the text continues normally.
The command does not insert an empty gap.

\subsubsection{Required Space}

The argument is a vertical length. It can be expressed using
physical units, the current line spacing or the text height:

\begin{lstlisting}[style=latex]
% Alternative ways to specify the required space:
\Needspace{3cm}
\Needspace{5\baselineskip}
\Needspace{0.25\textheight}
\end{lstlisting}

Here, \verb|\baselineskip| is the distance between consecutive
text baselines. The last example requests one quarter of the
text area's height.

\subsubsection{Headings and Following Text}

Place the command before a heading to leave room for both the
heading and the beginning of its text:

\begin{lstlisting}[style=latex]
\Needspace{7\baselineskip}
\subsection{Results}

The following results describe the response of the
system under the selected operating conditions.
\end{lstlisting}

The requested space should include the heading, its vertical
spacing and the desired number of text lines.

\subsubsection{Short Lists}

Reserve space before the introductory sentence to avoid
separating it from the beginning of a list:

\begin{lstlisting}[style=latex]
\Needspace{8\baselineskip}
The main assumptions are:

\begin{itemize}
    \item Steady flow.
    \item Constant fluid density.
    \item Negligible viscous losses.
\end{itemize}
\end{lstlisting}

To keep the entire list together, allow enough space for all
items and the spacing around them. Increase the requested
length if the items occupy multiple lines.

\subsubsection{Command Variants}

The package provides an approximate command and two variants
that check the available space directly:

\begin{lstlisting}[style=latex]
% Approximate space reservation
\needspace{5\baselineskip}

% Exact space check
\Needspace{5\baselineskip}

% Exact space check; preserves bottom alignment
\Needspace*{5\baselineskip}
\end{lstlisting}

The lowercase command is faster but less precise.
Both capitalized variants should be used between paragraphs.

When a page is ended early, \verb|\needspace| and
\verb|\Needspace| leave unused space at the bottom.
The starred variant respects \verb|\flushbottom| when
that setting is active:

\begin{lstlisting}[style=latex]
% Document setting
\flushbottom

% Later, between paragraphs
\Needspace*{6\baselineskip}
The following material requires additional room.
\end{lstlisting}

With \verb|\flushbottom|, stretchable vertical spacing can
expand to align the bottom of the page. The star does not
enable \verb|\flushbottom| by itself.

\subsubsection{Practical Limitations}

The package checks the requested length, not the actual
height of the following content. A paragraph or list that
exceeds that length can still continue on another page.

Use the command in the main text flow, before the relevant
material. It does not control the placement of floating
figures or tables. In a two-column layout, an insufficient
amount of space causes a break to the next column.